\documentclass[11pt,a4paper,oneside]{article}

% --- Paquetes Esenciales ---
\usepackage[utf8]{inputenc}
\usepackage[T1]{fontenc}
\usepackage[spanish]{babel}
\usepackage{geometry}
\geometry{margin=2.5cm}
\usepackage{amsmath, amssymb, amsfonts}
\usepackage{graphicx}
\usepackage{xcolor}
\usepackage{titlesec}
\usepackage{booktabs}
\usepackage{hyperref}
\usepackage{tcolorbox}
\usepackage{tikz}
\usetikzlibrary{shapes.geometric, arrows, positioning, shadows}

% --- Estética de los Títulos ---
\titleformat{\section}{\large\bfseries\color{darkblue}}{}{0em}{\thesection.\ }
\definecolor{darkblue}{rgb}{0.0, 0.2, 0.4}

% --- Estilos de Diagramas (TikZ) ---
\tikzstyle{block} = [rectangle, draw, fill=blue!10, text width=6em, text centered, rounded corners, minimum height=3em, shadow]
\tikzstyle{line} = [draw, -latex', thick]
\tikzstyle{cloud} = [draw, ellipse, fill=red!10, node distance=3cm, minimum height=2em]

% --- Metadatos ---
\title{\textbf{Neuro-Arquitectura de los Sistemas Agénticos}\\ \Large De la Psiquiatría Clínica al Silicio de la IA: Un Framework de Rigurosidad Científica}
\author{Ingeniería de Agentes Cognitivos \\ \textit{Whitepaper v2.0}}
\date{Abril 2026}

\begin{document}

\maketitle

\begin{abstract}
Este documento presenta una síntesis teórica y técnica sobre la convergencia entre la neurociencia cognitiva, la psiquiatría clínica y la ingeniería de agentes basados en modelos de lenguaje (LLMs). Se propone que la arquitectura de un agente moderno no es meramente una secuencia de llamadas a APIs, sino una réplica funcional de las funciones ejecutivas del cerebro humano. Exploramos marcos como la Teoría del Espacio de Trabajo Global (GWT), los Sistemas de Aprendizaje Complementarios (CLS) y el Cifrado Predictivo para proponer un modelo de diseño de agentes de alta fiabilidad.
\end{abstract}

\tableofcontents

\newpage

\section{Introducción: La Herencia de la Agencia}
La palabra \textbf{"agencia"} en el contexto de la Inteligencia Artificial no es un neologismo técnico, sino una transferencia conceptual de la psicología social. Según \textbf{Albert Bandura (1989)}, la agencia humana es la capacidad de ejercer influencia intencional sobre el funcionamiento propio.

\begin{tcolorbox}[colback=blue!5,colframe=blue!75,title=La Tétrada de Bandura en IA]
Cualquier framework agéntico (ej. LangGraph, CrewAI) debe satisfacer cuatro propiedades fundamentales:
\begin{itemize}
    \item \textbf{Intencionalidad:} El \textit{System Prompt} define el propósito.
    \item \textbf{Previsión:} El \textit{Planning} permite anticipar estados futuros.
    \item \textbf{Autorreactividad:} El \textit{ReAct Loop} permite ajustar la acción según el feedback.
    \item \textbf{Autoreflexión:} El \textit{Critic Layer} evalúa la eficacia del acto.
\end{itemize}
\end{tcolorbox}

\section{Analogías Funcionales: La Prótesis Prefrontal}
En neurociencia, la \textbf{Corteza Prefrontal (CPF)} actúa como el "CEO del cerebro". En la ingeniería de agentes, el orquestador o la lógica de grafos asume este rol de control ejecutivo.

\subsection{El Circuito Hipocampo-Neocórtex (CLS)}
La teoría de los \textbf{Sistemas de Aprendizaje Complementarios (CLS)} explica cómo el cerebro aprende:
\begin{enumerate}
    \item \textbf{Hipocampo:} Aprendizaje rápido de episodios (RAG / Memoria Episódica).
    \item \textbf{Neocórtex:} Aprendizaje lento de esquemas generales (Pesos del LLM / Parametric Memory).
\end{enumerate}
Esta separación es la solución biológica al \textit{Stability-Plasticity Dilemma}, el mismo que resolvemos en IA mediante el uso de bases de datos vectoriales para no tener que re-entrenar el modelo constantemente.

\begin{center}
\begin{tikzpicture}[node distance = 2cm, auto]
    % Nodes
    \node [block] (input) {Input / Percepción};
    \node [block, below of=input] (exec) {Ejecutivo (CPF)};
    \node [block, right of=exec, node distance=4cm] (mem) {Memoria Episódica (RAG)};
    \node [block, left of=exec, node distance=4cm] (neocortex) {Memoria Paramétrica (LLM)};
    
    % Lines
    \path [line] (input) -- (exec);
    \path [line] (exec) -- (mem);
    \path [line] (exec) -- (neocortex);
    \path [line] (mem) -- (exec);
    \path [line] (neocortex) -- (exec);
\end{tikzpicture}
\\ \textit{Figura 1: Arquitectura simplificada de un Sistema de Aprendizaje Complementario en Agentes.}
\end{center}

\section{Marcos Teóricos Avanzados}

\subsection{Teoría del Espacio de Trabajo Global (GWT)}
Propuesta por \textbf{Bernard Baars}, sugiere que la conciencia emerge de una "pizarra central" donde módulos especializados compiten por transmitir información. En sistemas Multi-Agente, esta arquitectura (\textit{Blackboard pattern}) es superior al paso de mensajes punto a punto, ya que mantiene una coherencia global del plan.

\subsection{Cifrado Predictivo y Energía Libre}
Siguiendo a \textbf{Karl Friston}, los sistemas inteligentes minimizan la \textit{Variational Free Energy} (sorpresa). Un agente robusto no solo predice el siguiente token, sino que \textbf{predice el resultado de sus acciones}. La divergencia entre la predicción y el resultado de una herramienta es el "Error de Predicción" que debe disparar la actualización del modelo interno del agente.

\section{Psicopatología y Diagnóstico de Sistemas}
Podemos diagnosticar fallos de agentes usando taxonomía psiquiátrica clínica:

\begin{table}[h]
\centering
\begin{tabular}{@{}lll@{}}
\toprule
\textbf{Patología Clínica} & \textbf{Manifestación en IA} & \textbf{Causa Técnica} \\ \midrule
Confabulación & Alucinaciones coherentes & Fallo en la verificación metacognitiva. \\
Fallo de Eferencia & Loops de repetición & El agente no reconoce su propia acción. \\
Estado Agéntico & Obediencia ciega dañina & Falta de guardrails morales internos. \\
Disfunción Ejecutiva & Pérdida de contexto & Ventana de contexto saturada o corrupta. \\ \bottomrule
\end{tabular}
\caption{Mapeo de patologías psiquiátricas a fallos de ingeniería agéntica.}
\end{table}

\section{Implementación: El Caso del \texttt{cognitive-code-agent}}
El análisis de la arquitectura del proyecto revela una implementación avanzada de estos principios:
\begin{itemize}
    \item \textbf{L0 (Working Memory):} Manejo de ventana de contexto con resumen dinámico.
    \item \textbf{L1 (Episodic Memory):} Implementación en Redis para persistencia de 90 días.
    \item \textbf{L2 (Semantic Memory):} Almacenamiento permanente en Milvus (Findings Store).
\end{itemize}

\section{Conclusión}
La ingeniería de agentes está evolucionando desde el "prompting" hacia la **Arquitectura Cognitiva**. El éxito de un AI Engineer especialista en agentes dependerá de su capacidad para diseñar sistemas que imiten no solo el lenguaje, sino los procesos biológicos de planificación, memoria y auto-monitoreo.

\end{document}
